% ============================================================
% PRESENTACIÓN DE LA MATERIA
% ============================================================
\section*{PRESENTACIÓN DE LA MATERIA}
\subsection*{Instrumentos Integrados}

\begin{center}
\small
\begin{tabular}{|p{0.32\linewidth}|p{0.58\linewidth}|}
\hline
\textbf{Escuela} & E.E.S.T. N° 7 ``Taller Regional Quilmes'' (IMPA) \\ \hline
\textbf{Especialidad} & Técnico/a en Aviónica \\ \hline
\textbf{Año} & 7mo año \\ \hline
\textbf{Campo} & Formación Técnico Específica \\ \hline
\textbf{Módulo} & Sistemas de Instrumentos de la Aeronave (Diseño Curricular 2025) \\ \hline
\textbf{Entorno} & Instrumentos Integrados -- Instrumentos Digitales (PFD / MFD) \\ \hline
\textbf{Carga horaria del módulo} & 4 horas semanales -- 144 horas anuales \\ \hline
\textbf{Marco curricular} & Diseño Curricular 2025 de la Educación Secundaria Técnica, Provincia de
Buenos Aires \\ \hline
\textbf{Docente} & Téc. Aviónico Julián Scicolone -- scicolonejulian@impatrq.com \\ \hline
\end{tabular}
\end{center}

% ------------------------------------------------------------
\section*{FUNDAMENTACIÓN}

El Diseño Curricular plantea que, en la atención de las aeronaves, es necesario conocer las
técnicas de \textbf{operación, montaje y mantenimiento de los sistemas de instrumentos
eléctricos, neumáticos, hidráulicos, de navegación e integrados}, y el uso de herramientas de mano
y especiales en prácticas de montaje y desmontaje de sus componentes. Estos sistemas establecen
un \textbf{cuarto nivel de complejidad} en las prácticas de mantenimiento: montaje y desmontaje,
inspección, evaluación del estado general de los componentes electrónicos y la tecnología de los
materiales. Se suman las capacidades de interpretación, organización, planificación y
certificación de los procesos, para trabajar en talleres de aviónica con su documentación, su
herramental específico y sus normas.

Este manual se centra en los \textbf{instrumentos integrados}: las pantallas de vuelo primarias
(\textbf{PFD}) y multifunción (\textbf{MFD}) que reemplazaron a los instrumentos de aguja, desde
los sistemas completos como el \textbf{Garmin G1000} y el \textbf{G5000} hasta los equipos de
reemplazo que hoy se instalan en la aviación general (\textbf{G5, GI 275, G3X Touch, TXi, AXIS,
Aspen}). Y en el oficio que hace posible instalarlos: \textbf{cables, terminales, crimpado,
conectores militares y D-sub, herramientas y ferretería} de la aviónica.

\begin{recuadro}
\textbf{La realidad del trabajo.} En un taller de aviónica de aviación general, gran parte del
trabajo diario es \textbf{modernizar paneles}: sacar instrumentos de vacío y de aguja, instalar
pantallas digitales, armar mazos de cables, crimpar cientos de contactos, configurar los equipos y
probarlos. Un mazo bien hecho funciona durante décadas; un contacto mal crimpado genera una falla
intermitente que nadie encuentra. Esta materia es ese oficio.
\end{recuadro}

% ------------------------------------------------------------
\section*{OBJETIVOS DE APRENDIZAJE}
Según el Diseño Curricular 2025, se espera que las y los estudiantes puedan:
\begin{itemize}
  \item Interpretar el funcionamiento de los componentes electrónicos y de los equipos, sobre la
        base de los principios físicos que regulan su funcionamiento y de la tecnología aplicada.
  \item Planificar y realizar el mantenimiento de los distintos equipos según las normativas
        vigentes y el criterio personal y profesional necesario.
  \item Determinar los métodos y los tiempos necesarios para las tareas de mantenimiento básico.
  \item Confeccionar y completar la documentación técnica de todas las etapas del mantenimiento.
  \item Desenvolverse profesionalmente dentro de un taller de aviónica, priorizando el trabajo
        grupal e interdisciplinario, la seguridad y el seguimiento de las normas vigentes.
  \item Interiorizarse en los avances y desarrollos de las nuevas tecnologías y la innovación
        digital, y su aplicación en el campo.
\end{itemize}

% ------------------------------------------------------------
\section*{ORGANIZACIÓN DEL MANUAL}

\begin{center}
\small
\begin{tabular}{|c|>{\raggedright\arraybackslash}p{0.52\linewidth}|>{\raggedright\arraybackslash}p{0.29\linewidth}|}
\hline
\textbf{Unidad} & \textbf{Contenidos} & \textbf{Contenido del Diseño Curricular} \\ \hline
1 & Del panel analógico a la cabina de cristal: T básica, EFIS, arquitecturas, normas y
certificación. & Clasificación y agrupación de instrumentos; sistemas electrónicos de vuelo \\
\hline
2 & Sensores: pitot-estático, ADC, AHRS, magnetómetro, sensores de motor; buses de datos. & Sistemas
de presión, ADC; actitud, rumbo; indicadores de motor \\ \hline
3 & PFD y MFD: presentación, alertas, banderas, modos reversionarios, instrumentos de reserva. &
Sistemas electrónicos de vuelo y de aviso \\ \hline
4 & Sistemas integrados Garmin G1000 / G1000 NXi y G5000. & Sistemas electrónicos de vuelo \\ \hline
5 & Equipos de reemplazo: G5, GI 275, G3X Touch, G500/G600 TXi, AXIS, Aspen. & Sistemas
electrónicos de vuelo; nuevas tecnologías \\ \hline
6 & Cables aeronáuticos, mazos, blindajes, masas y EWIS. & Instalación y normas: cables \\ \hline
7 & Terminales, empalmes, crimpado y soldadura. & Instalación y normas: pines \\ \hline
8 & Conectores circulares militares y D-sub. & Instalación y normas: pines, cables \\ \hline
9 & Herramientas del taller y ferretería eléctrica y electrónica aplicada. & Herramientas de mano y
especiales \\ \hline
10 & Instalación, configuración y documentación de un panel digital. & Instalación y normas
(tableros, iluminación) \\ \hline
11 & Ensayos, calibración, sistema pitot-estático y diagnóstico de fallas. & Ensayo, mediciones y
calibración; chequeo funcional \\ \hline
TP & Trabajos prácticos de taller. & Prácticas formativas \\ \hline
PI & Proyecto integrador: panel didáctico con su mazo de cables. & Prácticas formativas \\ \hline
\end{tabular}
\end{center}

\paragraph*{Metodología.} Cada unidad combina teoría, documentación real (manuales de instalación
y de mantenimiento de los fabricantes), ejemplos numéricos y actividades. Los \textbf{trabajos
prácticos} se hacen en el banco: crimpado, armado de conectores, mazos, ensayos de continuidad,
configuración y prueba pitot-estática. Cierra con un \textbf{proyecto integrador}: un panel
didáctico completo con su mazo de cables y su documentación.

\paragraph*{Fuentes.} Los datos de cada equipo (números de parte, conexiones, configuraciones) se
toman siempre del \textbf{manual de instalación vigente del fabricante}. Este manual da ejemplos
para aprender; no reemplaza esa documentación.

% ------------------------------------------------------------
\section*{SEGURIDAD EN EL TALLER DE AVIÓNICA}
\begin{recuadro}
\begin{enumerate}
  \item \textbf{Desenergizar} antes de conectar o desconectar equipos: batería desconectada o
        disyuntores abiertos con collar y tarjeta.
  \item \textbf{Descarga electrostática:} pulsera y superficie antiestática para manipular
        pantallas y computadoras; guardarlas en bolsas antiestáticas.
  \item \textbf{Herramientas de corte y pelado:} cuchillas siempre guardadas; nunca pelar cables
        hacia el cuerpo.
  \item \textbf{Pistola de calor y soldador:} superficies calientes; no dejarlos apoyados sobre
        cables ni sobre la mesa.
  \item \textbf{Sistema pitot-estático:} nunca soplar ni aplicar presión o vacío bruscos: se dañan
        los instrumentos y los sensores.
  \item \textbf{Orden y limpieza:} recortes de cable, contactos sueltos y virutas son
        \textbf{FOD}: pueden terminar dentro de un conector o de un equipo.
\end{enumerate}
\end{recuadro}

% ------------------------------------------------------------
\clearpage
\section*{HOJA DE RUTA DEL ESTUDIANTE}
\subsection*{Lo que vas a tener hecho al terminar el año}

\vspace{2pt}
{\small
\renewcommand{\arraystretch}{1.2}
\begin{tabular}{|c|>{\raggedright\arraybackslash}p{0.54\linewidth}|>{\centering\arraybackslash}p{0.08\linewidth}|>{\centering\arraybackslash}p{0.17\linewidth}|}
\hline
\textbf{OK} & \textbf{Logro} & \textbf{Unidad} & \textbf{Firma docente} \\ \hline
\casilla & Identificación de la T básica y de los datos de un PFD y un MFD & 1--3 & \\ \hline
\casilla & Diagrama de bloques de un sistema integrado (sensores, pantallas, buses) & 2--4 & \\
\hline
\casilla & Operación de un PFD y un MFD en simulador o entrenador & 3--4 & \\ \hline
\casilla & Cuadro comparativo de equipos de reemplazo (G5, GI 275, G3X, TXi, AXIS, Aspen) & 5 & \\
\hline
\casilla & Cálculo de sección de cable y caída de tensión & 6 & \\ \hline
\casilla & Crimpado de terminales y empalmes con ensayo de tracción & 7 & \\ \hline
\casilla & Terminación de un blindaje con manguito de soldadura & 7 & \\ \hline
\casilla & Armado de un conector D-sub de alta densidad con contactos crimpados & 8 & \\ \hline
\casilla & Inserción y extracción de contactos en un conector circular MIL-DTL-38999 & 8 & \\ \hline
\casilla & Inventario y uso correcto de las herramientas del taller & 9 & \\ \hline
\casilla & Plano de cableado de una instalación a partir del manual del fabricante & 10 & \\ \hline
\casilla & Mazo de cables armado, identificado y verificado pin a pin & 10 & \\ \hline
\casilla & Configuración y calibración de un equipo (o en modo entrenamiento) & 10 & \\ \hline
\casilla & Ensayo de pérdidas del sistema pitot-estático con banco de prueba & 11 & \\ \hline
\casilla & Diagnóstico de una falla simulada y tarjeta de trabajo completa & 11 & \\ \hline
\casilla & Proyecto integrador: panel didáctico presentado & PI & \\ \hline
\end{tabular}
}

\vspace{8pt}
\noindent Estudiante: \dotfill \quad Curso: \makebox[3cm]{\dotfill}

\vspace{6pt}
\noindent Logros completados: \makebox[1.5cm]{\dotfill} de 16 \hfill Fecha: \makebox[4cm]{\dotfill}
